\section{Background and Proof Obligations}
\subsection{Notation and elementary operations}
Write $[n,k,d]$ for a $k$-dimensional binary linear code of length $n$ with minimum nonzero Hamming weight exactly $d$, and $[n,k,\ge d]$ when only a lower bound is imposed. Let $\dopt(n,k)$ be the maximum achievable distance. Zero and repeated columns are allowed throughout. A code has \emph{full support} if no coordinate is identically zero. It is \emph{even} if all its words have even weight.

Puncturing deletes coordinates. Deleting $t<d(C)$ coordinates is injective on $C$ and reduces distance by at most $t$. Shortening on a coordinate set $S$ first restricts to words zero on $S$ and then deletes $S$. If the projection onto $S$ has rank $r$, shortening reduces dimension by exactly $r$ and does not reduce minimum distance. Parity extension sends $x$ to $(x,\sum_i x_i)$ and preserves dimension. These operations apply to all codes under consideration, without a symmetry or projectivity assumption.

\subsection{Classification and computational bounds}
Classical coding theory combines general inequalities with classification near extremal parameters. The Griesmer bound gives
\begin{equation}\label{eq:griesmer}
 n\ge g_2(k,d):=\sum_{j=0}^{k-1}\left\lceil d/2^j\right\rceil
\end{equation}
for a binary linear $[n,k,\ge d]$ code \citep{griesmer1960}. For the targets considered here, the decisive information is more specific than this bound. The eight-dimensional classification of \citet{bouyukliev2000} excludes exact $[41,8,18]$. We call this input \textbf{E1}. Together with $g_2(8,19)=42$, E1 also excludes $[41,8,\ge18]$.

\citet{bouyukliev2001} classify optimal binary linear codes of dimension at most seven. The input used here, \textbf{E2}, is an upper bound of 172 coordinate-equivalence classes of full-support $[40,7,18]$ codes. Our concrete catalogue contains 172 pairwise inequivalent representatives, so E2 establishes its completeness. The distribution by hull dimension is corroborated by \citet{bouyuklieva2025}: dimensions $1,3,5$ occur $33,62,77$ times. Hull counts are a consistency check; completeness follows from the class bound and actual inequivalence, not from matching these three numbers.

A separate classification technique extends residual codes. \citet{bouyuklieva2020}, Proposition~2, prove that a binary linear $[41,14,14]$ code does not exist. We call this input \textbf{E3}. Our fifth result imports the proposition as a whole and supplies the reduction from $[42,16,\ge13]$. The historical classification computations underlying E1--E3 remain published inputs; our computations concern the new extension obstruction and the explicit lower-bound matrices.

\subsection{Constructions, certification, and formal proof}
Code tables organize lower constructions and upper bounds at fixed parameters \citep{grasslTables}. They are useful starting points for identifying targets and recovering witnesses. A table entry, however, plays a different role from an explicit construction: the latter can be checked by row reduction and enumeration of all nonzero messages. We include the actual lower-bound generators in the appendix, including matrices reconstructed from the quasicyclic constructions recorded by CodeTables.

Computational nonexistence requires a finite domain with a justified connection to every hypothetical object. In the main proof, the classification fixes the mothers and a linear quotient fixes the coset domain. Two complete enumerators address coverage; a separate witness checker addresses rejection validity. These tasks are logically different. The rejected pairs have small witnesses, but the completeness of the accepted-coset lists still depends on exhaustive enumeration.

Lean~4 provides a further level of verification for definitions and deductions \citep{lean2021}; the formal developments use mathlib \citep{mathlib2020}. The formal statements refer to actual finite-dimensional subspaces, weights, and generator matrices. Their external hypotheses identify the precise facts supplied by published classification and computational evidence. This arrangement exposes where formal proof ends and ordinary computational trust begins, while checking the entire intervening mathematical argument.
